\section{基于共现计数的词向量方法}

\subsection{从``预测''到``计数''：另一种思路}

Word2Vec 是一种\textbf{基于预测}的方法：通过神经网络预测上下文词来学习词向量。但还有一种更直观的思路：既然我们关心的是词与词之间的共现关系，为什么不直接\textbf{统计}它们？

\begin{figure}[H]
\centering
\includegraphics[width=0.95\textwidth]{slides-images/slide-015.jpg}
\caption{为什么不直接统计共现次数？两种构建共现矩阵的方式：基于窗口和基于文档\protect\footnotemark}
\end{figure}
\footnotetext{Slides 第15页。}

构建共现矩阵 $X$ 有两种方式：
\begin{itemize}
    \item \textbf{基于窗口}：类似 Word2Vec，统计每个词在固定窗口内与其他词的共现次数，捕捉\textbf{语法和语义}信息
    \item \textbf{基于文档}：统计词在不同文档中的共现，捕捉\textbf{主题}信息（即潜在语义分析 LSA 的思路）
\end{itemize}

\subsection{共现矩阵示例}

\begin{figure}[H]
\centering
\includegraphics[width=0.95\textwidth]{slides-images/slide-016.jpg}
\caption{基于窗口的共现矩阵示例：语料为三个短句，窗口大小为 1\protect\footnotemark}
\end{figure}
\footnotetext{Slides 第16页。}

以一个极小的语料为例（``I like deep learning. I like NLP. I enjoy flying.''），窗口大小为 1，可以构建出如图所示的共现矩阵。矩阵中的每个元素 $X_{ij}$ 表示词 $i$ 和词 $j$ 在窗口内共现的次数。

\begin{warningbox}{共现矩阵的维度问题}
共现矩阵的大小为 $|V| \times |V|$（词表大小的平方）。如果词表有 40 万词，矩阵大小为 $400{,}000 \times 400{,}000$，远大于 Word2Vec 的 $400{,}000 \times d$（$d$ 通常为 100-300）。因此，直接使用共现矩阵作为词向量既\textbf{存储浪费}，又\textbf{稀疏低效}。
\end{warningbox}

\subsection{SVD 降维}

解决高维问题的经典方法是\textbf{奇异值分解}（Singular Value Decomposition, SVD）：

$$X = U \Sigma V^T$$

\begin{itemize}
    \item $U$ 和 $V$ 是正交矩阵（列向量互相正交且模为 1）
    \item $\Sigma$ 是对角矩阵，对角元素（奇异值）按大小递减排列
\end{itemize}

通过只保留最大的 $k$ 个奇异值（将其余的置零），我们得到矩阵 $X$ 的最优 $k$ 维近似，从而将高维共现矩阵压缩为低维词向量。

\begin{knowledgebox}{从 LSA 到现代词向量}
这种``共现矩阵 + SVD 降维''的方法最早被称为\textbf{潜在语义分析}（Latent Semantic Analysis, LSA），在心理学和信息检索领域有广泛应用。但直接对原始计数做 SVD 效果不佳。Doug Rohde 在 2000 年代的博士论文中发现了一系列改进技巧：对频率取对数、加权窗口距离、使用 Pearson 相关代替原始计数等。他甚至独立发现了线性语义成分的性质，但当时未受到广泛关注。
\end{knowledgebox}

\subsection{GloVe：统一计数与预测的模型}

\begin{figure}[H]
\centering
\includegraphics[width=0.95\textwidth]{slides-images/slide-020.jpg}
\caption{GloVe 的核心直觉：共现概率的\textbf{比率}能编码语义成分\protect\footnotemark}
\end{figure}
\footnotetext{Slides 第21页。Pennington, Socher, and Manning, EMNLP 2014。}

GloVe（Global Vectors for Word Representation）由 Stanford 的 Jeffrey Pennington、Richard Socher 和 Chris Manning 于 2014 年提出。它的核心洞察是：

\begin{importantbox}{GloVe 的关键直觉：共现概率的比率}
单独看一个词与 ice 或 steam 的共现概率，无法提取出有意义的语义成分。但如果看它们的\textbf{比率}：
\begin{itemize}
    \item $P(\text{solid} | \text{ice}) / P(\text{solid} | \text{steam})$ $\gg$ 1（solid 与 ice 相关）
    \item $P(\text{gas} | \text{ice}) / P(\text{gas} | \text{steam})$ $\ll$ 1（gas 与 steam 相关）
    \item $P(\text{water} | \text{ice}) / P(\text{water} | \text{steam})$ $\approx$ 1（water 对两者无偏）
\end{itemize}
这些比率\textbf{编码了 ice 与 steam 之间的``固态 vs 气态''语义差异}。
\end{importantbox}

为了将这种比率关系转化为向量空间中的\textbf{线性}关系，GloVe 引入了对数：比率的对数变为差值。因此 GloVe 的目标是学习词向量 $w_i$ 和 $\tilde{w}_j$，使得：

$$w_i^T \tilde{w}_j + b_i + \tilde{b}_j \approx \log X_{ij}$$

\begin{itemize}
    \item $X_{ij}$：词 $i$ 和词 $j$ 的共现次数
    \item $b_i, \tilde{b}_j$：偏置项
\end{itemize}

这样，两个词向量的差就对应于它们共现概率比率的对数，从而将语义差异编码为向量空间中的线性关系。

\begin{knowledgebox}{GloVe 与 Word2Vec 的关系}
表面上看，Word2Vec 是``基于预测''的方法，GloVe 是``基于计数''的方法。但实际上两者\textbf{深度相关}：
\begin{itemize}
    \item Word2Vec 的 Skip-gram 目标函数可以证明等价于对共现概率矩阵的隐式分解
    \item GloVe 的优化目标可以看作对共现矩阵的显式分解
    \item 两者在实践中产生质量相当的词向量
\end{itemize}
区别更多在于实现细节和训练效率，而非根本原理。
\end{knowledgebox}

\subsection{本章小结}

基于共现计数的方法（SVD、LSA、GloVe）提供了学习词向量的另一条路径。GloVe 通过建立共现概率比率与向量差之间的对应关系，将``计数''与``预测''两种方法的优势统一起来。其关键洞察是：共现概率的\textbf{比率}（而非绝对值）编码了有意义的语义成分，而\textbf{取对数}将比率转化为向量空间中的线性关系。

%% ========== Part 4: 词向量的评估 ==========

